% ___________________________________________________________________________

\hypertarget{secuencias}{}
\pdfbookmark[0]{Secuencias}{secuencias}

\begin{frame}[containsverbatim]
\frametitle{Secuencias}
\small
\begin{itemize}
\item {} 
Son argumentos dedicados a contar.

\item {} 
Uso:
\begin{itemize}
\item contador en bucles
\item definir rangos de valores
\item crear vectores fila.
\end{itemize}

\end{itemize}

Ejemplo:
\begin{Octave}
  >> 1:5 
  ans =
    1 2 3 4 5
  >> 0:2:10
  ans =
    0 2 4 6 8 10
  >> 0:-3:-10
  ans =
    0 -3 -6 -9
\end{Octave}

\end{frame}


\begin{frame}[containsverbatim]
\frametitle{Nuestro primer script desde qtOctave}

\begin{itemize}

\item Editar y grabar en disco un fichero llamado
    ``tablavalores.m'' con el siguiente contenido:

\begin{Octave}[morekeywords={x,y}]
  x=-1:0.1:1; 
  y=1./(1+x.*x);
\end{Octave}

\item Los operadores \verb|.*| y \verb|./| multiplican y dividen
  elemento a elemento.

\item {} En la ventana de comandos, ejecutar el script y pedir los valores
  de \emph{x} e \emph{y}

\begin{Octave}[morekeywords={tablavalores,x,y}]
  >> tablavalores 
  >> x 
    (...valores de x...)  
  >> y
    (...valores de y...)
\end{Octave}


\item Otras funciones disponibles disponibles en la ventana de comandos:
  \texttt{ls}, \texttt{pwd}, \texttt{cd}
\end{itemize}





\end{frame}


% % ___________________________________________________________________________

%\subsection{Matrices}
% ___________________________________________________________________________

\hypertarget{tipos-num-ricos}{}
\pdfbookmark[0]{Tipos numéricos}{tipos-num-ricos}
\begin{frame}[containsverbatim]
\frametitle{Tipos numéricos}

\begin{itemize}
\item 
  \emph{Por omisión}, Octave (al igual que Matlab) emplea un único
  tipo numérico: 
\begin{center}
  \Large
  \textbf{\alert{reales de doble precisión}}
  \\[0.5em]
  \normalsize
  (arrays n-dimensionales de)
  \\[0.4em]~
\end{center}


\item Cualquier otro tipo tendrá que asignarse explícitamente:
\begin{flushright}
  \verb|int8|, \verb|int16|, \verb|int32|, \verb|int64|
  \\
  \small \verb|uint8|, \verb|uint16|, \verb|uint32|, \verb|uint64|
  \\
  \verb|single|, \verb|double|
\end{flushright}

\item Ejemplo: \verb|i=int8(32)|
\end{itemize}
\end{frame}



% ___________________________________________________________________________

\hypertarget{para-diferenciar-filas-y-columnas}{}
\pdfbookmark[1]{Para diferenciar filas y columnas}{para-diferenciar-filas-y-columnas}
\begin{frame}[containsverbatim]
\frametitle{Para escribir matrices}

\textbf{Notación:}
\begin{itemize}
\item Se utilizan corchetes
\item El espacio en blanco o la coma separan columnas en una fila
\item El retorno de carro o el punto y coma separan filas
\end{itemize}

\textbf{Ejemplo} 
\begin{Octave}
>> A=[1,2; 3,4]
A =
   1   2
   3   4
\end{Octave}

\textbf{Ejercicio:} Escribir de dos formas distintas la matriz:
$$
M=\left(
\begin{array}{lll}
  11 & 12 & 13 \\
  21 & 22 & 23
\end{array}
\right)
$$
\end{frame}

% ___________________________________________________________________________

\hypertarget{sub-ndices}{}
\pdfbookmark[0]{Subíndices}{sub-ndices}
\begin{frame}[containsverbatim]\frametitle{Subíndices}

  Los arrays son indexables del modo usual (uso de paréntesis, \texttt{1} es el primer índice)
\begin{Octave}
  >> M(1,2) 
  ans = 12
  >> M(end,end) 
  ans = 23
\end{Octave}

\small
\vfill\hfill ...por supuesto, \texttt{end} es el \emph{ultimo} índice.
\end{frame}

% ___________________________________________________________________________


\begin{frame}[containsverbatim]
\frametitle{El lenguaje de Matlab/Octave es muuuy flexible}
\begin{Octave}
  >> a = pi 
  a = 3.1416
  >> a(1) 
  ans = 3.1416
  >> a(1,1) 
  ans = 3.1416
  >> a(1,1,1) 
  ans = 3.1416
\end{Octave}
\begin{itemize}
\item {} 
Lo anterior devolvería un error en cualquier otro lenguaje de programación.

\item {} 
¿Esto es bueno o es malo?

\end{itemize}

\end{frame}

% ___________________________________________________________________________

\begin{frame}[containsverbatim]
\frametitle{Al menos, hay verificación de límites}
\begin{Octave}
  >> a=rand(3,3) 
  a = 
 
    0.583220   0.285967   0.010804 
    0.336715   0.530313   0.860391 
    0.060404   0.999098   0.501339 
 
    >> a(3,4) 
    error: invalid column index = 4
\end{Octave}
  
\end{frame}

% ___________________________________________________________________________

\hypertarget{pero-no-para-la-asignaci-n}{}
\pdfbookmark[0]{¡Pero no para la asignación!}{pero-no-para-la-asignaci-n}
\begin{frame}[containsverbatim]
\frametitle{¡Pero no para la asignación!}

\begin{Octave}
  >> nuevo(4)=1 
  nuevo = 
 
     0   0   0   1 
 
  >> nuevo(4,3)=2 
  nuevo = 
 
     0   0   0   1 
     0   0   0   0 
     0   0   0   0 
     0   0   2   0
\end{Octave}


\end{frame}

% ___________________________________________________________________________

\hypertarget{submatrices}{}
\pdfbookmark[0]{Submatrices}{submatrices}
\begin{frame}\frametitle{Submatrices}

Supongamos que de la matriz $M$ queremos extraer una
matriz $S$ formada por las 3 últimas filas y columnas:

$$
M=\left(\begin{array}{rrrrr}
  11 & 12 & 13 & 14 & 15 \\
  21 & 22 & 23 & 24 & 25 \\
  31 & 32 & \alert{33} & \alert{34} & \alert{35} \\
  41 & 42 & \alert{43} & \alert{44} & \alert{45} \\
  51 & 52 & \alert{53} & \alert{54} & \alert{55}
\end{array}\right)
$$
\end{frame}

% ___________________________________________________________________________

\begin{frame}[containsverbatim]
\frametitle{Submatrices}

Los subíndices se seleccionan mediante secuencias.  Por ejemplo, para
extraer la submatriz formada por las filas
\begin{itemize}
\item \emph{i} de 3 a 5
\item \emph{j} de 3 a 5
\end{itemize}
escribimos:
\vspace{-0.6em}
\begin{Octave}
>> S = M(3:5,3:5) 
 33   34   35 
 43   44   45 
 53   54   55
\end{Octave}

\medskip
Para extraer elementos de una fila o columna, podemos usar
'\verb|:|'. Por ejemplo:
\vspace{-0.5em}
\begin{Octave}
>> S = M(3:5,:) 
 31   32   33   34   35 
 41   42   43   44   45 
 51   52   53   54   55
\end{Octave}


\end{frame}

% ___________________________________________________________________________


\begin{frame}[containsverbatim]
\frametitle{Submatrices}

\textbf{Ejercicio}: La orden \verb|pascal(N)| devuelve una matriz
cuadrada de orden $N$ formada a partir del respectivo triángulo de
Pascal. 
\begin{itemize}
\item Calcular la matriz de Pascal $P$ de orden $6$.
\item 
Extraer la matriz $T$ formada por:
\begin{itemize}
\item Las filas $1$ a $3$
\item Las columnas $2$ y $5$
\end{itemize}
\end{itemize}
\end{frame}

\begin{frame}[containsverbatim]
\frametitle{Submatrices}

\textbf{Solución}: La matriz formada por las
\begin{itemize}
\item filas $1$ a $3$ y...
\item columnas $2$ y $5$
\end{itemize}
se obtiene de la siguiente forma:

\begin{Octave}
>> P = pascal(6)
>> T = P(1:3,[2,5]);
\end{Octave}

\end{frame}

% ___________________________________________________________________________

\begin{frame}[containsverbatim]
\frametitle{Submatrices}
\textbf{Ejercicio}: Ahora queremos extraer la matriz formada por las
columnas $1$, $3$ y $5$.
\hfill
\end{frame}

\begin{frame}[containsverbatim]
\frametitle{Submatrices}
\textbf{Ejercicio}: Ahora queremos extraer la matriz formada por las
columnas $1$, $3$ y $5$.

\hfill

\textbf{Solución}: La matriz formada por las
columnas $1$, $3$ y $5$ es...

\begin{Octave}
>> P(:,1:2:5) 
\end{Octave}
 
\end{frame}

% ___________________________________________________________________________

\hypertarget{operadores}{}
\pdfbookmark[0]{Operadores}{operadores}
\begin{frame}[containsverbatim]
\frametitle{Operadores}
\begin{description}[itemsep=1.1em]
\item [Operadores matriciales:]
\verb|+    -    *    /    ^|

\item [Operadores escalares:]
\verb|.*    ./    .^|

\item [Operadores lógicos:]
\verb# &   |    ! #

\item [Operadores de comparación:]
\verb|<    >    ==    >=    <=    !=|

\item [Operadores matriciales o de conjuntos:]
\verb# &&   ||#

\end{description}

\end{frame}

% ___________________________________________________________________________

\begin{frame}[containsverbatim]
\frametitle{Diferencias entre \& y \&\&}
\begin{Octave}
>> a=[1,2;0,1]; 
>> b=[1,0;0,1]; 
>> a&b 
ans = 
 
   1   0 
   0   1 
 
>> a&&b 
ans = 0
\end{Octave}

\begin{itemize}
\item 
  \verb|&&| se utiliza para las sentencias condicionales.
\item 
  \verb|&&| es un ``operador de cortocircuito''
\item 
  \textbf{Ejercicio}: ¿Cuándo es \verb'A && A' verdadero? 
%  ¿y \verb' A || A'? ¿y \verb'A && ~A'? ¿y \verb'A || ~A'?
  La orden \verb'false && 1/0' ¿devuelve un error (división por cero)?
\end{itemize}
\end{frame}

% ___________________________________________________________________________

\hypertarget{operadores-ii}{}
\pdfbookmark[0]{Operadores II}{operadores-ii}
\begin{frame}[containsverbatim]
\frametitle{Operadores II}

Operadores matriciales y escalares pueden confundirse fácilmente:
\begin{Octave}
>> a=rand(3,3);b=rand(3,3); 
>> a*b 
ans = 
    1.0297    0.9105    0.3293 
    0.9663    0.8267    0.4211 
    0.5355    0.4318    0.3279 
>> a.*b 
ans = 
    0.1824    0.3253    0.0563 
    0.5500    0.6003    0.1897 
    0.0458    0.0017    0.1822
\end{Octave}
  
\end{frame}

% ___________________________________________________________________________

\hypertarget{operadores-iii}{}
\pdfbookmark[0]{Operadores III}{operadores-iii}
\begin{frame}[containsverbatim]
\frametitle{¿Cuál es la solución ``correcta''?}

\begin{Octave}
>> a=[1,2,3;4,5,6;7,8,9]; 
>> a.^pi 
ans = 
    1.0000    8.8250   31.5443 
   77.8802  156.9925  278.3776 
  451.8079  687.2913  995.0416 
>> a^pi 
ans = 
 1.0e+03 * 
 0.6943-0.0004i 0.8540-0.0001i 1.0136+0.0002i 
 1.5743-0.0000i 1.9344-0.0000i 2.2946+0.0000i 
 2.4543+0.0003i 3.0149+0.0001i 3.5756-0.0002i
\end{Octave}

\end{frame}

% ___________________________________________________________________________

\begin{frame}\frametitle{Álgebra}

Con estas funciones se puede crear casi cualquier matriz
\begin{description}
\item [eye:]
Matriz de ceros con unos en la diagonal


\item [linspace:]
Vector de elementos equiespaciados


\item [logspace:]
Vector de elementos espaciados exponencialmente


\item [ones:]
Matriz de unos

\item [diag:]
Matriz con la diagonal que decidamos

\item [zeros:]
Matriz de ceros


\item [rand:]
Matriz de números pseudoaleatorios.


\item [meshgrid:]
Matrices equiespaciadas de dos dimensiones

\end{description}

\end{frame}

% ___________________________________________________________________________

\hypertarget{manipulaci-n-de-matrices}{}
\pdfbookmark[0]{Manipulación de matrices}{manipulaci-n-de-matrices}
\begin{frame}\frametitle{Álgebra}
Manipulación de matrices
\begin{description}
\item [reshape:]
Cambia la forma de la matriz conservando el número de
elementos


\item [transpose:]
Traspuesta. Equivale a .'


\item [ctranspose:]
Matriz conjugada. Equivale a '


\item [rot90:]
Gira la matriz 90 grados en sentido antihorario.


\end{description}
Tamaño de una matriz
\begin{description}
\item[size]  Devuelve el número de filas y columnas
\item[length] Devuelve el número de elementos de un vector\footnote{En
    realidad, \texttt{length(A)=max(size(A))}}
\end{description}

\end{frame}

% ___________________________________________________________________________

\begin{frame}[containsverbatim]
\frametitle{Resolución de sistemas lineales}

Para resolver sistemas de eucuaciones lineales contamos con un
operador universal:
\begin{Octave}
>> A=[1,0;2,1]; y=[2;4]; 
>> x=A\y 
x = 
  2 
  0
\end{Octave}
\begin{itemize}\item 
  A bajo nivel se usan las (excelentes) bibliotecas ATLAS y LAPACK
\end{itemize}
\end{frame}

%%% Local Variables: 
%%% mode: latex
%%% TeX-master: "modulo_octave"
%%% End: 
